I explained that autism is one way a person's brain can work and that autistic people are not all the same. Communication may include speech, gestures, writing, pictures, AAC, movement, or silence. Some students may need quiet, space, clear directions, processing time, predictable routines, or a break from noise.

Using the OAR friendship guidance, I taught five actions:
\begin{enumerate}
\item \textbf{Invite} without forcing participation.
\item \textbf{Wait} so the person has time to process, communicate, or try.
\item \textbf{Ask} respectfully instead of guessing what someone needs.
\item \textbf{Respect} sensory needs, personal space, communication tools, breaks, and privacy.
\item \textbf{Include} students who are new, quiet, tired, or participating differently.
\end{enumerate}

The practice scenario came from placement patterns: during a coding activity, one student already knew how to enter the program, several classmates were still stuck, and another learner was new to the group. Students suggested showing one step without grabbing the Chromebook, using a calm voice, waiting for the learner to try, and treating the new learner as part of the class. The key distinction was between helping and taking over.

I closed with: \textit{One way I can be a supportive classmate is... because...} The word ``because'' required students to connect an action to reducing stress, increasing access, protecting dignity, or helping someone belong.